\chapter{Richard Ewen Borcherds}\label{ch1998richardborcherds:richard-ewen-borcherds-fields-medal-1998-a-comprehensive-report-on-his-work-and-impact}

\begin{center}\rule{0.5\linewidth}{0.5pt}\end{center}

\textbf{Prepared:} July 2026
\textbf{Subject:} Richard Ewen Borcherds (born January 29, 1959, Oxford, England)
\textbf{Position:} Professor of Mathematics, University of California, Berkeley (since 1993)
\textbf{Doctoral Advisor:} John H. Conway (University of Cambridge)
\textbf{Fields Medal Citation:} \emph{``For his contributions to algebra, in
particular the theory of vertex algebras, the Moonshine conjecture,
and the Monster group.''}

\begin{center}\rule{0.5\linewidth}{0.5pt}\end{center}

\begin{prereqs}
Group theory (sporadic groups, representation theory),
Lie algebras, modular forms, combinatorics.
\end{prereqs}

\begin{objectives}
State the Monstrous Moonshine conjecture and its proof; explain the
Monster group and the $j$-function connection; describe vertex
algebras and their role.
\end{objectives}

% ============================================================
\section{Problem History}
% ============================================================

\textbf{Monstrous Moonshine} refers to the mysterious connection
between the Monster group $\mathbb{M}$---the largest sporadic
finite simple group, of order approximately $8\times 10^{53}$---
and modular functions, particularly the $j$-function.  The
connection was observed by Conway and Norton in 1979.

The \textbf{Monster group} $\mathbb{M}$ was predicted by Fischer
and Griess (1973) and constructed by Griess (1982) as
automorphisms of a 196,884-dimensional algebra.  The dimension
196,884 is tantalizingly close to 196,883, the first non-trivial
coefficient of the $j$-function:
\[
  j(q) = q^{-1} + 196884\,q + 21493760\,q^2 + \cdots.
\]

Conway and Norton formulated the \textbf{Moonshine conjecture}:
there exists a graded representation
$V=\bigoplus_{n=-1}^{\infty}V_n$ of $\mathbb{M}$ such that each
Hauptmodul for a genus 0 group agrees with the graded character of
$\mathbb{M}$.  Equivalently, the McKay--Thompson series:
\[
  T_g(q) = \sum_{n=-1}^{\infty}\mathrm{tr}(g|V_n)\,q^n
\]
is the $j$-function for $g=1$ and a Hauptmodul for each
non-identity element.

The \textbf{196,884 problem} had fascinated mathematicians since the
Monster's discovery.  Thompson and McKay noted
$196884 = 196883+1$, and the next coefficient
$21493760 = 21296876+196883+1$, where these are dimensions of
Monster representations.

Borcherds, while a student of Conway, began working on Moonshine in
the mid-1980s, developing \textbf{vertex algebras} as the natural
algebraic framework.

Frenkel and Lepowsky had introduced vertex algebras in 1984 in the
context of string theory, but the connection to the Monster was not
understood.

The term ``moonshine'' was chosen by Conway to convey the whimsical
nature of the connection---the idea that the Monster and modular
functions should be related seemed as unlikely as illicit
whiskey-making.  Yet the connection turned out to be profound.

A \textbf{Hauptmodul} (German for ``main module'') for a Kleinian
group $\Gamma$ of genus zero is a meromorphic function that
generates the field of meromorphic modular functions for $\Gamma$.
For the full modular group $\mathrm{SL}(2,\mathbb{Z})$, the
Hauptmodul is the $j$-function itself.  For congruence subgroups,
it is a function with a single pole per fundamental domain.

% ============================================================
\section{How Was It Solved}
% ============================================================

Borcherds's proof \cite{borcherds1992} proceeds in several stages.

\textbf{Stage 1: The Monster vertex algebra.}
He constructs the \textbf{Monster vertex algebra} $V^{\natural}$,
an infinite-dimensional graded algebra with a Lie bracket-like
structure.  Vertex algebra axioms encode locality of fields in 2D
conformal field theory: for any $a,b\in V$, vertex operators
$Y(a,z)$ and $Y(b,w)$ commute when $|z|<|w|$.

A \textbf{vertex algebra} $V$ consists of a vector space with a
state $|0\rangle$, translation operator $T$, and vertex operators
$Y(a,z) = \sum_{n\in\mathbb{Z}}a_n z^{-n-1}$ satisfying locality,
associativity, and the Jacobi identity.  The \textbf{griess
algebra} is the subspace $V_2$ of the Monster vertex algebra,
equipped with a commutative non-associative product
$a\circ b = a_1 b$.

The Monster vertex algebra is constructed from the Moonshine module
using the Griess algebra and Monster representation theory.  The
key property is that $V^{\natural}$ is a \textbf{holomorphic vertex
operator algebra} of central charge 24, meaning its 1-point
functions on the torus transform as modular forms.

\textbf{Stage 2: The No-ghost theorem.}
Borcherds proves the \textbf{No-ghost theorem}: the space of
physical states in a string-theoretic construction has no
negative-norm states.  This tool, adapted from Goddard--Thorn,
provides the mechanism for proving McKay--Thompson series are
Hauptmoduln.

\textbf{Stage 3: The denominator identity.}
Borcherds defines the \textbf{denominator identity} for the Monster
Lie algebra, an infinite-dimensional Lie algebra from the Monster
vertex algebra:
\[
  e^{\rho} - e^{-\rho}
  \prod_{\alpha>0}(1-e^{-\alpha})^{m(\alpha)}
  = j(q) - j(z)
\]
where $m(\alpha)$ are multiplicities.  This remarkable identity
connects the Monster and the $j$-function.  The Monster Lie algebra
is constructed from $V^{\natural}$ as a generalized Kac--Moody
algebra with root spaces indexed by a 2-dimensional lattice
$\mathrm{II}_{1,1}$.

\textbf{Stage 4: The Conway--Norton conjecture.}
He proves the \textbf{Conway--Norton conjecture} using the
denominator identity, modular invariance of the vertex algebra, and
the faithful Monster action on $V^{\natural}$.

The proof also introduces \textbf{generalized Kac--Moody algebras}
(Borcherds--Kac--Moody algebras), allowing imaginary simple roots.

% ============================================================
\section{Impact of This Discovery}
% ============================================================

The proof of Moonshine is one of the most remarkable achievements
in 20th-century mathematics, establishing deep connections between
finite group theory, modular functions, and string theory.

The Monster vertex algebra $V^{\natural}$ has become fundamental in
algebra and mathematical physics.  It is the unique holomorphic
vertex operator algebra of central charge 24 with trivial Griess
algebra.

Vertex algebras have become central tools in string theory,
conformal field theory, and representation theory of
infinite-dimensional Lie algebras, encoding the operator product
expansion in 2D conformal field theory.

Generalized Kac--Moody algebras have found applications in Lie
algebras, algebraic groups, and number theory, generalizing the
Weyl--Kac character formula.

Moonshine has been generalized to other sporadic groups:
\textbf{Mathieu moonshine}, \textbf{UMoonshine}, and
\textbf{Umbral moonshine}, connecting finite groups to modular forms
and mock modular forms.

The Monster vertex algebra has a remarkable property: it is the
unique holomorphic vertex operator algebra of central charge 24
with the Monster as automorphism group.  This uniqueness, proved by
Dong, Lin, and others, is analogous to the uniqueness of the Leech
lattice.

% ============================================================
\section{Modern Views and Improvements}
% ============================================================

Mathieu moonshine, connecting $M_{24}$ to the elliptic genus of K3
surfaces, has been proved by Gaberdiel, Hampapura, and Mukhi.

UMoonshine, proposed by Duncan, Griffin, and Ono, connects $M_{24}$
to mock modular forms, leading to the theory of Umbral moonshine.

The classification of holomorphic vertex operator algebras of
central charge 24, completed by Dong, Lin, and others, shows
$V^{\natural}$ is one of exactly 71 such algebras, connecting to
Niemeier lattices and the Leech lattice.

Recent developments include \textbf{K3 moonshine}, which connects
the elliptic genus of K3 surfaces to mock modular forms and
$M_{24}$ representations.

The \textbf{Umbral moonshine} conjecture, proved by Duncan, Griffin,
and Ono, establishes moonshine connections for all 23 Niemeier
lattices, not just the Leech lattice.

The \textbf{Mathieu moonshine} connection between $M_{24}$ and the
elliptic genus of K3 surfaces was discovered by Eguchi, Ooguri, and
Tachikawa in 2010.

Modern research explores connections to the Langlands program,
arithmetic geometry, and quantum gravity.  The Monster vertex
algebra appears in string theory on $T^{24}$.

Borcherds continues contributing to vertex algebras, Lie algebras,
and sporadic groups, including work on the fake monster Lie
superalgebra and generalized moonshine.

% ============================================================
\begin{center}\rule{0.5\linewidth}{0.5pt}\end{center}

\section*{Summary}
\addcontentsline{toc}{section}{Summary}

Borcherds proved the Moonshine conjecture by constructing the
Monster vertex algebra and proving McKay--Thompson series are
Hauptmoduln, using the No-ghost theorem, generalized Kac--Moody
algebras, and vertex operator algebras.  His work established a
profound connection between the Monster group, modular functions,
and string theory.

% ============================================================
